% Physical pages 3 and 4 of the concise study: the movement of the formulary,
% stated as an argument and not as a summary of what follows. The developed
% comparison of the three readings begins on the fifth page.

\section{The Propers: Themes and Movement}
\zlabel{triptych:concise:themes:start}

This Mass asks for a peace it does not define, shows a power exercised on
earth, and confesses in its single prayer that without God we cannot please
God. It opens with \latin{Da pacem}, a request, and a destination: ``We shall
go into the house of the Lord''. It closes with a command to bring sacrifices
and come into the Lord's courts, and a prayer to be made worthy of what has
been received. Between them stand a thanksgiving for a grace already given to a
divided church, a man carried to Jesus and forgiven before he is healed, and
Moses at an altar in the sight of Israel.

The Introit's antiphon is cited to Ecclus 36:18, the close of the son of
Sirach's prayer for ``Jerusalem, the city which thou hast sanctified''. The
chant keeps that verse's middle clause word for word, \latin{ut prophetae tui
fideles inveniantur}, and changes its frame: it asks \latin{pacem} where the
Vulgate asks \latin{mercedem} and the Greek a reward, adds \latin{Domine}, and
prays for \latin{servi tui, et plebis tuae Israel}, one servant and the people,
where the Vulgate has many servants. The words \latin{Da pacem} stand in no
verse of the Clementine Vulgate. The psalm verse is the opening of Ps~121, a
Song of Ascents, the pilgrim's joy at the call to go up to the house of the
Lord, and the same verse returns after the Epistle as the Gradual's respond. The
destination named at the entrance is named a second time.

The Collect asks that \latin{tuae miserationis operatio}, the working of God's
mercy, may direct our hearts, and gives as its reason \latin{quia tibi sine te
placere non possumus}. The one we are to please is God, and the one without
whom we cannot please him is God himself; the 1861 Philadelphia missal renders
\latin{sine te} ``without thy help''. The prayer is older than the Sunday
number the Missal gives it. The Old Gelasian Sacramentary sets it at the head of
a Mass headed only \latin{Item alia missa}, with this Sunday's Secret and
Postcommunion; the Gregorian sacramentary in H.~A. Wilson's edition sets the
three at the Nineteenth Sunday, and the 1962 Missal at the Eighteenth. It is
the Mass's only oration.

The Epistle is the thanksgiving that opens St Paul's first letter to Corinth.
He gives thanks ``for the grace of God that is given you in Christ Jesus'':
they are rich ``in all utterance and in all knowledge'', nothing is wanting to
them in any grace as they wait for the Lord's revelation, and he ``will confirm
you unto the end without crime''. The name of Christ stands in four of the five
verses. The lesson stops one verse before \latin{Fidelis Deus} and two before
the appeal ``that there be no schisms among you'' (1:10), so that the church
thanked for its riches is the church about to be told of its parties. The
lesson breaks into the Sunday lessons from Ephesians, which resume on the
Nineteenth Sunday, and the Sunday itself follows the September Ember Days. Bl.\
Ildefonso Schuster heads it \latin{Tertia post natale Sancti Cypriani}, the name
the seventh-century Roman gospel list of Würzburg already gives the Sunday of
this Gospel, and explains that at Rome the Ember Saturday's vigil Mass served
as the Sunday's own. In the Würzburg epistle list, which Germain Morin
published as the oldest lectionary of the Roman Church, the lesson already
stands with its present extent, under the feast of the angel and immediately
after the lessons of the September Ember Saturday.

The Gradual sings verses 1 and 7 of Ps~121. Between them the psalm sets the
pilgrims' feet in the courts of Jerusalem, ``built as a city, which is compact
together'', and bids them ``pray ye for the things that are for the peace of
Jerusalem''; verse 7 is that prayer, \latin{Fiat pax in virtute tua: et
abundantia in turribus tuis}. The Introit's peace now has a place: a city with
strength and towers. The Alleluia takes one verse of Ps~101, the prayer of the
poor man, at the point where its lament has turned to Sion: ``The Gentiles shall
fear thy name, O Lord, and all the kings of the earth thy glory.'' The Missal
drops the connective \latin{Et} and stops before the reason the next verse
gives, ``For the Lord hath built up Sion''. The chants that name the house are
old at this Sunday. In four of the six manuscripts of Hesbert's
\work{Antiphonale Missarum Sextuplex}, as the gregorien.info database reports
them, \latin{Laetatus sum} is already its Gradual, and the same four give its
Offertory and Communion, three of them its Introit; \latin{Timebunt gentes}
stands in all six, but as a Gradual at other Sundays.

The Gospel, Mt 9:1--8, stands between the Gerasenes who begged Jesus to leave
their coasts and the call of Matthew at the custom house. Jesus crosses the
lake ``and came into his own city''. Men bring a paralytic lying on a bed;
``seeing their faith'', Jesus says to the sick man, ``Be of good heart, son,
thy sins are forgiven thee''. Some of the scribes say within themselves, ``He
blasphemeth''. Jesus asks which is easier to say, and to show ``that the Son of
man hath power on earth to forgive sins'' he bids the man rise, take up his bed
and go home. The man does so, and the crowds fear and glorify God ``that gave
such power to men''. The Creed follows.

The Offertory takes the Mass to the foot of Sinai. Its antiphon is a
compilation and not a quotation of the verses it cites. Where Exodus says that
Moses built an altar, \latin{aedificavit}, the antiphon says that he hallowed
it, \latin{sanctificavit}, the verb the Vulgate uses of the altar's
consecration; it adds \latin{sacrificium vespertinum in odorem suavitatis}, the
language of the daily evening offering; and its last words, \latin{in conspectu
filiorum Israel}, stand in the same chapter at 24:17, where the glory of the
Lord is seen on the mountain. The chapter it cites goes on to the blood of the
covenant sprinkled on the altar and the people. The antiphon once had verses in
which Moses intercedes for the people after the golden calf; the 1962 Missal
sings the antiphon alone.

The Secret says what the altar is for: \latin{per huius sacrificii veneranda
commercia}, through the venerable exchanges of this sacrifice, God makes us
\latin{unius summae divinitatis participes}, sharers in the one supreme
Godhead. The verb, \latin{efficis}, is present. The prayer asks that as we know
God's truth, \latin{cognoscimus}, so we may attain it by a worthy life,
\latin{assequamur}; and the Missal directs the Preface of the Most Holy Trinity
after it. The Communion sings the middle of Ps~95:8--9, a psalm headed in the
Vulgate ``when the house was built after the captivity'': \latin{Tollite
hostias, et introite in atria eius}. Where the Vulgate has \latin{in atrio sancto
eius}, the Missal sings \latin{in aula sancta eius}, as the Pustet Missal of
1862 already does. The Postcommunion begins \latin{Gratias tibi referimus},
thanks for the sacred gift by which the communicants have been
\latin{vegetati}, quickened, and asks \latin{ut dignos nos eius participatione
perficias}. The Old Gelasian carries that prayer twice, in the Mass it shares
with the Collect and the Secret and again as the Postcommunion of a Mass for
the dedication of a place that had been a synagogue. The Mass whose Epistle begins \latin{Gratias
ago} ends its prayers with thanks.

Three habits run through the whole. The first is the house. The Introit and the
Gradual go up to it, the Alleluia sees the nations come to Sion as it is built,
the Communion enters its courts, and in the Gospel the forgiven man is sent
\latin{in domum tuam}. The second is the visible given for the invisible. The
Gospel proves an unseen forgiveness by a seen healing; the rest of the
formulary takes visible things, an altar, victims, an evening sacrifice, a gift
received, and claims an invisible work done through them, \latin{efficis} in
the Secret and \latin{vegetati} in the Postcommunion. The third is that nothing
in the Mass is self-made. The Collect says so outright; the Epistle's riches are
``given''; the paralytic is carried in, and says nothing; and the last prayer
asks God himself to make the communicants worthy of what they have already
received.

Three readings of the whole Mass follow from these texts, and the expansive
study develops each at length. The first hears the formulary from its chants,
as one ascent from the prayer for peace to the entry into the Lord's courts, and
finds the peace of the city in the charity of which it is built. The second
hears it from the Gospel, as the Son of man's power on earth to forgive, shown
in his own city and set beside the altar of the covenant whose blood, in the
words of institution, is shed ``for the remission of sins''. The third hears it
from the Collect, as a Mass that gives thanks at both ends for what no one in it
earned, with a man at its centre who was carried. They are complementary
emphases, three governing questions put to one set of texts; each accounts for
something the others leave aside, and at one point, the crowd's praise at the
end of the Gospel, two of them part over whose judgement of it to follow. The commentary that
follows takes them up together, at the points where they answer the same
question.
\zlabel{triptych:concise:themes:end}
